\section*{Chapter \ref{ch:b3:identical-particles} --- Identical Particles}

\begin{solution}{exo:b3:identical-particles:1}
(a) Swapping twice is doing nothing: \omterm{def:b3:quantum-formalism:observable}{eigenvalues} $\pm1$. (b)
\omterm{def:b3:identical-particles:exchange}{Identical particles} enter $\hat H$ symmetrically, so $\hat H\hat
P_{12} = \hat P_{12}\hat H$: the symmetry class never changes. (c)
$\ket a\otimes\ket b$ and $\ket b\otimes\ket a$ would be
\emph{different} states describing the same physical situation ---
an unobservable distinction the formalism must not carry. (d)
$^1$H (2 fermions): boson; deuterium (3): fermion; $^4$He: boson;
$^3$He: fermion; H$_2$ (4): boson; photon: boson.
\end{solution}

\begin{solution}{exo:b3:identical-particles:2}
(a) $(\ket a\otimes\ket b - \ket b\otimes\ket a)/\sqrt2$. (b)
Identically zero. (c) The alternating sum over the six
permutations --- the $3\times3$ determinant of orbitals against
particles; swapping two rows flips a determinant's sign. (d) A
determinant with two equal columns vanishes: two fermions in one
orbital is a repeated column.
\end{solution}

\begin{solution}{exo:b3:identical-particles:3}
(a) $3^2 = 9$. (b) $3$ doubles $+$ $3$ pairs $= 6$. (c) $3$ pairs.
(d) Six single-particle states (three orbitals $\times$ two spins):
$\binom62 = 15$.
\end{solution}

\begin{solution}{exo:b3:identical-particles:4}
(a) C: $1s^22s^22p^2$; N: $2p^3$; O: $2p^4$; Ne: $2p^6$. (b)
Nitrogen's three $2p$ electrons occupy the three orbitals singly,
spins parallel. (c) Parallel spins force an antisymmetric spatial
state, whose exchange hole spares Coulomb repulsion --- alignment
is bought with electrostatic savings. (d) Helium (with its noble
kin); fluorine, one electron short of neon's closure, tears one
from almost anything.
\end{solution}

\begin{solution}{exo:b3:identical-particles:5}
(a) Bosons: $2E_1$; spinless fermions: $E_1 + 4E_1 = 5E_1$;
electrons: $2E_1$ (opposite spins share $n = 1$). (b) Bosons and
electrons: $E_1 + 4E_1 = 5E_1$; spinless fermions: $E_1 + 9E_1 =
10E_1$. (c) Compressing fermions forces occupation of ever-higher
levels: the energy cost per squeeze is an outward pressure that no
cooling removes --- matter's incompressibility in embryo. (d)
$\sum_{n=1}^N n^2 \approx N^3/3$: the Fermi sea's total grows as
$N^3$.
\end{solution}

\begin{solution}{exo:b3:identical-particles:6}
(a) The two terms cancel at $\vect r_1 = \vect r_2$. (b) Expanding,
the cross terms contribute $\mp2|\bra a\hat{\vect r}\ket b|^2$:
the antisymmetric state has the \emph{larger} mean square
separation. (c) Farther apart means less repulsion energy: $E_- <
E_+$. (d) Helium's triplet (antisymmetric space) lies below the
same configuration's singlet --- orthohelium under parahelium, by
\qty{0.8}{eV} of spared repulsion.
\end{solution}

\begin{solution}{exo:b3:identical-particles:7}
(a) $1s^2$ is a symmetric spatial state: the spins must form the
antisymmetric singlet. (b) Splitting $= 2J$: $J \approx
\qty{0.4}{eV}$. The genuine magnetic dipole--dipole energy at
$a_0$: $\mu_0\mu_{\text{B}}^2/4\pi a_0^3 \approx \qty{4e-4}{eV}$
--- a thousand times smaller: the ``spin force'' is electrostatics
wearing spin's colours. (c) Decay to $1\,^1S$ needs a spin flip
(singlet--triplet) \emph{and} an $s \to s$ transition ($\Delta l =
0$): doubly forbidden. (d) $8000$ s against \qty{1.6}{ns}: twelve
orders of magnitude, from two selection rules --- metastability is
rules stacked, not forces weakened.
\end{solution}

\begin{solution}{exo:b3:identical-particles:8}
(a) Total exchange sign $=$ (spin part) $\times (-1)^J$ must be
$-1$: triplet (symmetric spin) takes odd $J$, singlet even $J$.
(b) $3:1$ from the multiplicities. (c) The 3/4 ortho fraction
relaxes to para ($J = 0$), releasing $2B \approx \qty{15}{meV}$
per molecule --- \emph{more} than the latent heat of the liquid:
an uncatalysed tank slowly boils itself dry. Plants convert
ortho $\to$ para over a catalyst during liquefaction. (d) Liquid
hydrogen rocketry --- every launch pad inherits a nuclear-spin
statistics problem.
\end{solution}

\begin{solution}{exo:b3:identical-particles:9}
(a) Rates go as the amplitude squared: $n + 1$ into an occupied
mode, $n$ for the reverse. (b) $10^{10} + 1$ against $1$: the
crowded mode wins by ten billion. (c) One spontaneous photon in
the favoured mode tilts the odds; each stimulated photon raises
$n$ and steepens the tilt: an avalanche of identical photons ---
gain narrowing into a single mode. (d) In $\hat a^\dagger\ket n =
\sqrt{n+1}\,\ket{n+1}$: the bosonic mode is the oscillator, and
the enhancement is its ladder coefficient.
\end{solution}

\begin{solution}{exo:b3:identical-particles:10}
(a) Each gas expands into double the volume: $\Delta S =
Nk_{\text{B}}\ln2$ per gas. (b) Closing and reopening the tap
would manufacture entropy without any change of state ---
perpetual bookkeeping fraud. (c) With named molecules, the
combined gas has $\binom{2N}{N}$ extra ``which-side'' assignments
worth $2Nk_{\text{B}}\ln2$; dividing all state counts by $N!$ ---
the honest count for \omterm{def:b3:identical-particles:exchange}{identical particles} --- removes exactly this.
(d) The $N!$ that classical physics needed and could not justify
is the symmetrisation postulate seen from afar.
\end{solution}

\begin{solution}{exo:b3:identical-particles:11}
(a) $E_{\text{F}} = \hbar^2(3\pi^2n_{\text{e}})^{2/3}/2m_{\text{e}}
\approx \qty{0.16}{MeV}$ --- a third of the electron's \omterm{def:b3:relativistic-dynamics:momentum-energy}{rest
energy}: the star sits at the edge of the relativistic regime. (b)
$k_{\text{B}}T \approx \qty{0.9}{keV} \ll E_{\text{F}}$: on the
Fermi scale the white-hot star is deeply degenerate --- ``cold''.
(c) The electrons stack Pauli-fashion up to
$\qty{0.16}{MeV}$; compressing the star raises every occupied
level, and that energy gradient is a pressure independent of
temperature --- exclusion, not heat, carries the weight. (d)
Relativistic degeneracy softens the pressure law until gravity
wins at a finite mass: the Chandrasekhar limit, kept for
\cref{ch:b3:astrophysics}.
\end{solution}

\begin{solution}{exo:b3:identical-particles:12}
(a) $\hat{\vect S}_1\cdot\hat{\vect S}_2 = \tfrac12(\hat S^2 -
\tfrac32\hbar^2)$: $-\tfrac34\hbar^2$ on the singlet,
$+\tfrac14\hbar^2$ on the triplet --- the stated \omterm{def:b3:hamiltonian-mechanics:hamiltonian}{Hamiltonian}
reproduces the $2J$ split. (b) A lattice of aligned neighbours:
spontaneous magnetisation --- ferromagnetism. (c) $J/k_{\text{B}}
\approx \qty{1200}{K}$: iron's \qty{1043}{K} is exactly this
scale. (d) Dipole--dipole magnetism is $10^{-4}$ eV: it would
order at millikelvins. Ferromagnetism is Coulomb repulsion
choosing spin arrangements through exchange geometry.
\end{solution}

\begin{solution}{pb:b3:identical-particles:1}
\textbf{1.} Orbitals $(n, l, m)$ with $n^2$ per $n$; spin doubles;
Pauli caps one electron per full label: $2$, $8$, $18$.
\textbf{2.} Inner electrons screen the nucleus differently for
different $l$: penetrating $s$ orbitals sink below less
penetrating $p$ and $d$ --- energy order follows $(n, l)$, not $n$
alone.
\textbf{3.} Order $1s\,2s\,2p\,3s\,3p\,4s\,3d$; Na: [Ne]$3s^1$;
Cl: [Ne]$3s^23p^5$; K: [Ar]$4s^1$; Fe: [Ar]$4s^23d^6$.
\textbf{4.} Rows close at each noble configuration: $2$ ($1s$),
$8$ ($2s2p$), $8$ again ($3s3p$ --- because $3d$ lies above $4s$
and waits), then $18$ ($4s3d4p$).
\textbf{5.} A closed shell: large gap to the next orbital, no
cheap electron to give, no room to take.
\textbf{6.} $Z_{\text{eff}}$ is the nuclear charge minus the
screening of inner (and partially, same-shell) electrons: sodium's
valence electron sees $11 - 10 \approx 1$; helium's two share the
bare $Z = 2$, each screening the other by only a third.
\textbf{7.} Li $\to$ Be rises, drops at B; N $\to$ O drops:
the two dips.
\textbf{8.} Boron opens $2p$, higher and less penetrating than
$2s$: its new electron is simply easier to remove.
\textbf{9.} Oxygen must, for the first time, \emph{pair} two
electrons in one $p$ orbital: extra repulsion, lost exchange
stabilisation --- the fourth $2p$ electron is discounted.
\textbf{10.} Each electron screens the other imperfectly:
the fit $24.6 = 13.6\,Z_{\text{eff}}^2$ gives $Z_{\text{eff}} =
1.34$ --- two-thirds of a charge screened.
\textbf{11.} A nearly complete core screen ($Z_{\text{eff}}
\approx 1.3$) and the outer electron's $n = 2$: both cut the
binding.
\textbf{12.} Added electrons enter the \emph{same} shell and
screen each other poorly while $Z$ climbs: $Z_{\text{eff}}$ grows
and the whole shell contracts --- neon is smaller than lithium.
\textbf{13.} One cheap $s$ electron over a noble core: ionisation
$\sim\qty{5}{eV}$, instant $+1$ ions, violent electron donation
--- the family traits are one configuration.
\textbf{14.} One hole in a $p^6$ closure: huge electron
affinity, $-1$ ions, and pairing of holes into X$_2$ molecules.
\textbf{15.} $2s$ and $2p$ lie close: promotion $2s^22p^2 \to$
four half-filled hybrids costs little and buys two extra bonds ---
carbon tetravalence is a near-degeneracy exploited.
\textbf{16.} With four electrons, both the bonding \emph{and}
antibonding combinations fill: the gains cancel, no covalent
bond --- helium remains a loner.
\textbf{17.} The $d$ electrons bury themselves inside the $4s$
skin: outer chemistry, set by the $s$ shell, changes little across
ten elements --- a chemical plateau.
\textbf{18.} Scandium, [Ar]$4s^23d^1$: the first unpaired $d$
electron.
\textbf{19.} A magnetic moment of three aligned spins
($\sim3\mu_{\text{B}}$); the alignment is paid by exchange ---
Coulomb savings, not magnetic forces.
\textbf{20.} Slightly distinguishable electrons could all crowd
toward $1s$ (no exact cancellation of symmetrised states): shells
would leak, closures blur, and the periodic law dissolve.
\textbf{21.} $v/c \approx Z\alpha = 0.58$: relativistic mass
increase contracts and stabilises the $6s$ orbital; gold's
$5d \to 6s$ absorption slides into the blue (leaving reflected
yellow), and mercury's tightened $6s^2$ pair bonds so feebly the
metal is liquid.
\textbf{22.} They run nearly level: occupation-dependent
ordering, $4s$ ionising first, and the $3d^5$/$3d^{10}$ borrowings
of Cr and Cu --- half-filled and filled subshells buy exchange
stability cheaply.
\textbf{23.} Chemistry is the electrons' affair, and the
electrons feel only the nuclear \emph{charge}: mass enters merely
through tiny vibrational (isotope) shifts.
\textbf{24.} Shell structure (favouring noble-gas logic) against
relativistic stabilisations and sheer Coulomb strain at
$Z\alpha \to 1$: superheavy chemistry is their tug-of-war.
\textbf{25.} A catalogue of orbitals, one doubling, one minus
sign: rows $2, 8, 8, 18$, dips at boron and oxygen, families from
alkali to noble, the ionisation ladder from \qty{5.1}{eV} to
\qty{24.6}{eV} --- the periodic law derived, not memorised.
\end{solution}
